\chapter{Discussion}
\label{chap:discussion}
\lhead{\emph{Discussion}}

This chapter places the results in context, highlights practical advantages, and outlines limitations and threats to validity.

% -----------------------------------------------------------------------------
\section{Advantages of the Proposed System}
\label{sec:discussion-advantages}

\subsection{Explainability as a First-Class Output}

Most systems add explanations after the fact. This system builds explanations into the data model. A \texttt{ThreatAssessment} cannot exist without a hypothesis and evidence list. A \texttt{DecisionPlan} includes per-action rationale. A \texttt{PolicyDecision} records the rule that approved, held, or denied each action. This makes the audit trail a direct record of what the system actually used at decision time, not a reconstruction. \cite{nist80061}

\subsection{Safety Independent of Model Quality}

The policy engine enforces hard constraints even when the model is wrong. Actions that lack required entities or violate thresholds are denied before they can reach execution. This is stronger than typical playbook systems where safety depends on the correctness of the playbook. \cite{tilbury2024humans}

\subsection{Operational Feasibility}

The custom MLP planner runs in under 5 ms. The TF-IDF scorer is sub-millisecond. This keeps the total intelligence overhead under 10 ms per alert, which is practical for real-time triage. The seq2seq alternative is far slower, which makes it difficult to deploy at scale. \cite{ban2023alertfatigue}

\subsection{Model-Agnostic Architecture}

The model registry separates model choice from orchestration code. Profiles can be switched without redeploying the engine. This makes experimentation and rollback safe and fast.

% -----------------------------------------------------------------------------
\section{Comparison with Traditional SIEM and SOAR Approaches}
\label{sec:discussion-comparison}

\subsection{Against SIEM Platforms}

SIEM platforms collect and correlate logs, but they do not reason about response. \cite{nist80092} They surface alerts and rely on analysts to decide what to do next. This system starts where SIEM stops. It takes normalized alerts and adds scoring, planning, policy checks, and audited execution.

\subsection{Against Playbook-Based SOAR}

SOAR platforms execute predefined workflows. That works well for known cases, but it struggles with new or mixed patterns. \cite{tilbury2024humans} The planner in this system learns from data and generalizes across scenarios. It still makes mistakes, but it does not fail silently when no playbook matches.

\subsection{Against Generative LLM-Based Planning}

LLM planning can generate plausible output but often hallucinates entity values. In a response system, a wrong IP or host is a serious safety problem. The structured prediction approach avoids this by design: the model selects action types, and parameters are pulled from the alert context. This keeps actions grounded in real entities.

% -----------------------------------------------------------------------------
\section{Limitations}
\label{sec:discussion-limitations}

\subsection{Evaluation on Synthetic Data Only}

All quantitative results are based on synthetic data generated from the same pipeline used for training. This makes the evaluation controlled but optimistic. Real SOC data will be messier and more diverse. \cite{sommer2010outside}

\subsection{Limited Alert Type Coverage}

Only five alert types were used in training. Commercial SIEMs cover hundreds. Expanding coverage is possible but requires additional data and retraining.

\subsection{Simulated SIEM Connector}

The system has a real Wazuh connector, but most evaluation runs use the simulator. A production integration would require careful mapping and testing against real SIEM payloads.

\subsection{Manually Authored Policy Rules}

Policy rules are configured manually. This requires domain expertise and tuning. There is no automated policy learning in the current implementation.

\subsection{Feedback Loop Not Fully Closed}

Feedback is captured in audit logs, but online weight updates are not integrated into the live inference loop. Retraining is still an offline step.

% -----------------------------------------------------------------------------
\section{Threats to Validity}
\label{sec:discussion-threats}

\subsection{Internal Validity}

Training and test data are generated by the same pipeline. This can inflate accuracy because the model learns generator patterns rather than real-world structure.

\subsection{External Validity}

Generalization to real SOC data is not proven. Distribution shift is a known challenge for security ML systems. \cite{sommer2010outside}

\subsection{Construct Validity}

Strategy accuracy and action F1 measure label agreement, not operational impact. Whether the system reduces analyst workload or improves response time is not measured.

\subsection{Statistical Validity}

Results are reported from single training runs with fixed seeds. Multiple seeds and confidence intervals would strengthen the claims, but were not included in this work.
